\section{Logger Utility}
\label{sec:logger_utility}

\subsection{Overview}
The \texttt{Logger} class in \texttt{MLGWorks.Utils.Logging} provides a flexible and centralized logging framework for Unity applications. It captures Unity logs, routes them into categorized log files, and supports features like file cleanup, multiple output targets, and per-level filtering.

\subsection{LogLevel and Targets}
Each log entry has a severity level:
\begin{itemize}
  \item \texttt{Debug}
  \item \texttt{Info}
  \item \texttt{Warning}
  \item \texttt{Error}
\end{itemize}

Log levels are directed to specific files using 5-bit masks:
\begin{itemize}
  \item Bit 0: Debug file
  \item Bit 1: Info file
  \item Bit 2: Warning file
  \item Bit 3: Error file
  \item Bit 4: Combined file (receives all)
\end{itemize}

Example:
\begin{verbatim}
public int debugTargets = (1 << 0) | (1 << 4); // debug + combined
\end{verbatim}

\subsection{Configuration Options}

\begin{description}
  \item[\texttt{LogLocationType}] Determines the base path used for log storage:
    \begin{itemize}
        \item \texttt{PersistentDataPath}
        \item \texttt{DataPath}
        \item \texttt{Custom}
    \end{itemize}

  \item[\texttt{relativePath}] Folder path relative to base directory.
  \item[\texttt{customPath}] Absolute path if \texttt{Custom} type is selected.
  \item[\texttt{fileExtension}] Default: \texttt{.log}.
  \item[\texttt{maxLogFileCount}] Limits how many timestamped log file sets are retained. Use \texttt{-1} to keep all.
  \item[\texttt{enableDebugLogging}] Toggles whether \texttt{Debug()} logs are allowed.
\end{description}

\subsection{Logging API}

The logging methods enqueue messages for writing in a background-safe fashion:
\begin{itemize}
  \item \texttt{Logger.Debug(string msg)}
  \item \texttt{Logger.Info(string msg)}
  \item \texttt{Logger.Warning(string msg)}
  \item \texttt{Logger.Error(string msg)}
\end{itemize}

\begin{verbatim}
Logger.Info("App started.");
\end{verbatim}

All logs are timestamped and formatted like:
\begin{verbatim}
[14:23:01.123] [Info] App started.
\end{verbatim}

\subsection{Unity Integration}

Unity messages (e.g., \texttt{Debug.Log}, \texttt{Debug.LogError}) are automatically captured via \texttt{Application.logMessageReceivedThreaded} and redirected to the logger with a \texttt{[UNITY]} prefix.

\subsection{Lifecycle}

\begin{itemize}
  \item Logging begins lazily when the first log is emitted.
  \item \texttt{Flush()} is called on every frame.
  \item Log files are closed on \texttt{OnDestroy()} and \texttt{OnApplicationQuit()}.
\end{itemize}

\subsection{Log Cleanup}

The method \texttt{CleanupOldLogFiles()} deletes old log groups based on timestamps:
\begin{itemize}
  \item Files are grouped by prefix timestamp (e.g., \texttt{2025-07-20\_12-30-00}).
  \item Groups exceeding \texttt{maxLogFileCount} are deleted oldest-first.
\end{itemize}

\subsection{Example Usage}
\begin{verbatim}
Logger.Debug("This is a debug log.");
Logger.Info("App launched successfully.");
Logger.Warning("Disk space running low.");
Logger.Error("Unhandled exception occurred.");
\end{verbatim}

\subsection{Testing}
\texttt{EmitTestLogs()} can be called to quickly generate log messages of all levels.

\subsection{Editor Support}
In the Unity Editor, \texttt{FlushAndShutdown()} can be used to immediately close files and clear writers without waiting for app shutdown.

\subsection{Implementation Notes}
\begin{itemize}
  \item Uses a thread-safe \texttt{ConcurrentQueue} to stage logs.
  \item Writes are flushed in \texttt{Update()}.
  \item Log file creation uses timestamps to separate sessions.
  \item Each log level writes to target files based on its configured mask.
\end{itemize}
